\documentclass[aspectratio=169,10pt]{beamer}

\usepackage{mathtools}
\usepackage[scaled=.98,sups]{XCharter}
\usepackage[charter,vvarbb]{newtxmath}
\PassOptionsToPackage{nomail,rule}{beamerthemeAmurmaple}
\PassOptionsToPackage{table}{xcolor}
\usepackage{beamerthemeAmurmaple}
\usepackage{colortbl}
\usepackage{microtype}
\usepackage{booktabs}
\usepackage{graphicx}
\usepackage{ragged2e}
\usepackage{array}
\usepackage{listings}
\usepackage{tikz}
\usetikzlibrary{arrows.meta,positioning,calc,matrix,fit,backgrounds,decorations.pathreplacing,shapes.multipart}
\graphicspath{{assets/generated/}{assets/}{../tikz/output/}{../figures/output/}}

\definecolor{Ink}{HTML}{202A34}
\definecolor{SoftGray}{HTML}{667580}
\definecolor{CodeBg}{HTML}{F4F5F6}
\definecolor{CodeBlue}{HTML}{005A9C}
\definecolor{CodeGreen}{HTML}{287A4B}
\definecolor{WarmOrange}{HTML}{DF6C00}
\colorlet{Accent}{AmurmapleRed}
\colorlet{AccentWash}{AmurmapleRed!8}

\setbeamercolor{normal text}{fg=Ink,bg=white}
\setbeamerfont{footnote}{size=\fontsize{7}{8}}
\setbeameroption{hide notes}
\hfuzz=6pt

\lstdefinestyle{lecturecode}{
  basicstyle=\ttfamily\fontsize{6.75}{7.85}\selectfont,
  keywordstyle=\bfseries\color{CodeBlue},
  commentstyle=\itshape\color{CodeGreen},
  stringstyle=\color{WarmOrange!80!black},
  backgroundcolor=\color{CodeBg},
  frame=single,
  rulecolor=\color{SoftGray!35},
  numbers=left,
  numberstyle=\tiny\color{SoftGray},
  numbersep=4pt,
  xleftmargin=1.2em,
  framexleftmargin=1.0em,
  showstringspaces=false,
  columns=fullflexible,
  keepspaces=true,
  breaklines=true,
  tabsize=2,
  morekeywords={concept,requires,constexpr,decltype,nullptr,noexcept,static_assert,std,size_t,int32_t,int64_t,uint32_t,uint64_t,template,typename,class,struct,delete,default,explicit,auto}
}

\newcommand{\nnz}{\operatorname{nnz}}
\newcommand{\flops}{\operatorname{flops}}
\newcommand{\fmt}[1]{\textcolor{Accent}{\bfseries #1}}
\newcommand{\keyterm}[1]{\textcolor{Accent}{\bfseries #1}}
\newcommand{\metric}[1]{\colorbox{AccentWash}{\strut\textcolor{Ink}{\bfseries #1}}}
\newcommand{\metricbox}[1]{%
  \par\vspace{0.25em}%
  \centering%
  \colorbox{AccentWash}{\parbox{0.94\linewidth}{\centering\strut\textcolor{Ink}{\bfseries #1}}}%
  \par\vspace{0.10em}%
}
\newcommand{\src}[1]{\par\vspace{0.12em}{\fontsize{6.3}{7.4}\selectfont\color{SoftGray}#1\par}}

\newif\ifprogressive
\progressivefalse
\newcommand{\step}{\ifprogressive\pause\fi}
\newcommand{\sectionfoot}[1]{\gdef\insertsectionhead{#1}}
\newcommand{\subhead}[1]{%
  \medskip{\large\bfseries\textcolor{Accent}{#1}\par}%
  \vspace{-0.18cm}{\color{Accent}\rule{0.88\linewidth}{1.2pt}\par}\smallskip}

\tikzset{
  flow/.style={draw=Accent,rounded corners=2pt,fill=Accent!7,minimum height=0.72cm,align=center,font=\small},
  data/.style={draw=WarmOrange,rounded corners=2pt,fill=WarmOrange!9,minimum height=0.72cm,align=center,font=\small},
  codebox/.style={draw=CodeBlue,rounded corners=2pt,fill=CodeBg,font=\ttfamily\scriptsize,align=left},
  arrow/.style={-{Latex[length=2.0mm]},thick,SoftGray}
}

%%%%%%%%%%%%%%%%%%%%%%%%%%%%%%%%%%%%%%%%
%% Title Slide
%%%%%%%%%%%%%%%%%%%%%%%%%%%%%%%%%%%%%%%%
\title[C++ Metaprogramming in Sparse HPC]{C++ Metaprogramming \& Zero-Cost Abstractions \\ in Sparse Matrix Computation}
\subtitle{E399 / E599 - Large Scale Sparse Computation, Fall 2026}
\author[Yuxi Hong]{Yuxi Hong}
\position{Assistant Professor}
\event{E399 / E599}
\institute[IU Luddy ISE]{%
  Department of Intelligent Systems Engineering\\
  Luddy School of Informatics, Computing, and Engineering\\
  Indiana University, Bloomington%
}
\titlegraphic{\includegraphics[height=0.98cm,keepaspectratio]{assets/iu_trident.png}}

\begin{document}

\begin{frame}[plain]
  \titlepage
  \note{[Lecture Objectives]\par
    1. Understand the zero-overhead principle in high-performance computing.\par
    2. Master 3-tier type parameterization (IT, NT, OT) to minimize memory bandwidth.\par
    3. Learn compile-time safety: move semantics, deleted copy operations, and concepts.\par
    4. Explore advanced techniques: constexpr static dispatch and expression templates.}
\end{frame}

\agendaframe[Agenda][Metaprogramming in Sparse HPC]{
  \agendaitem{01}{Motivation: The Zero-Overhead Principle in Sparse HPC}
  \agendaitem{02}{3-Tier Type Parameterization: Decoupling IT, NT, and OT}
  \agendaitem{03}{Resource Ownership, RAII, and Compile-Time Safety}
  \agendaitem{04}{C++20 Concepts \& Expressive Compile-Time Diagnostics}
  \agendaitem{05}{Static Dispatch \& Branch Elimination with \texttt{if constexpr}}
  \agendaitem{06}{Expression Templates \& Single-Pass Kernel Fusion}
  \agendaitem{07}{Policy-Based Design for Extensible Sparse Frameworks}
}

%%%%%%%%%%%%%%%%%%%%%%%%%%%%%%%%%%%%%%%%%%%%%%%%%%%%%%%%%%%%%%%%%%%%%%%%%%%%%%%
\section{01: Motivation: The Zero-Overhead Principle in Sparse HPC}
\sectionfoot{01: Zero-Overhead Abstraction}
%%%%%%%%%%%%%%%%%%%%%%%%%%%%%%%%%%%%%%%%%%%%%%%%%%%%%%%%%%%%%%%%%%%%%%%%%%%%%%%

\begin{frame}[fragile]{The Cost of Traditional Abstractions in HPC}
  \begin{columns}[T,onlytextwidth]
    \begin{column}{0.48\textwidth}
      \subhead{Why OO Polymorphism Fails in Sparse HPC}
      \begin{itemize}\setlength{\itemsep}{0.45em}
        \item \keyterm{Memory-Bound Kernel}: SpMV has very low arithmetic intensity ($\text{FLOPs/Byte} \ll 1$). Every unnecessary byte transferred hurts performance.
        \item \keyterm{vtable Indirection}: Virtual function calls prevent compiler inlining and disable loop SIMD vectorization.
        \item \keyterm{Type Erasure Overhead}: Using \texttt{void*} or base pointers destroys compile-time type knowledge and alignment info.
      \end{itemize}
    \end{column}
    \begin{column}{0.48\textwidth}
      \subhead{The C++ Zero-Overhead Promise}
      \begin{block}{\large Bjarne Stroustrup's Rule}
        \textit{"What you don't use, you don't pay for. And further: What you do use, you couldn't hand code any better."}
      \end{block}
      \vspace{0.5em}
      \begin{itemize}\setlength{\itemsep}{0.4em}
        \item \textbf{Compile-Time Specialization}: Emit machine instructions tailored exactly to the precision, matrix sparsity, and memory layout.
        \item \textbf{Zero Runtime Latency}: Eliminate dynamic branches and allocations from the critical path.
      \end{itemize}
    \end{column}
  \end{columns}
\end{frame}

%%%%%%%%%%%%%%%%%%%%%%%%%%%%%%%%%%%%%%%%%%%%%%%%%%%%%%%%%%%%%%%%%%%%%%%%%%%%%%%
\section{02: 3-Tier Type Parameterization: Decoupling IT, NT, and OT}
\sectionfoot{02: 3-Tier Type Parameterization}
%%%%%%%%%%%%%%%%%%%%%%%%%%%%%%%%%%%%%%%%%%%%%%%%%%%%%%%%%%%%%%%%%%%%%%%%%%%%%%%

\begin{frame}[fragile]{Why Decouple Index, Value, and Offset Types?}
  \subhead{The Big-Graph Dilemma: $m = 50\text{M}$ Vertices, $\nnz = 3\text{B}$ Edges}
  \vspace{-0.2cm}
  \begin{itemize}
    \item \textbf{Vertices ($m$)}: $50 \times 10^6 < 2^{31}-1 \implies$ Fits comfortably in 32-bit integer (\texttt{int32\_t}, 4 bytes).
    \item \textbf{Edges ($\nnz$)}: $3 \times 10^9 > 2^{31}-1 \implies$ Exceeds 32-bit limit, requires 64-bit integer (\texttt{int64\_t}, 8 bytes).
  \end{itemize}

  \begin{table}
    \centering\small
    \begin{tabular}{lcccc}
      \toprule
      \textbf{Design Pattern} & \textbf{Index (\texttt{IT})} & \textbf{Offset (\texttt{OT})} & \textbf{Values (\texttt{NT})} & \textbf{Total Memory Footprint} \\
      \midrule
      Monolithic 32-bit & \texttt{int32\_t} & \texttt{int32\_t} & \texttt{double} & \textbf{OVERFLOW / CRASH} ($\nnz > 2^{31}$) \\
      Monolithic 64-bit & \texttt{int64\_t} & \texttt{int64\_t} & \texttt{double} & $3\times 8 + 3\times 8 = \mathbf{45.08\text{ GB}}$ \\
      \rowcolor{AccentWash}
      \textbf{SPCraft 3-Tier Decoupled} & \textbf{\texttt{int32\_t}} & \textbf{\texttt{int64\_t}} & \textbf{\texttt{double}} & $3\times 4 + 3\times 8 = \mathbf{33.90\text{ GB}}$ (\fmt{24.8\% Saved!}) \\
      \bottomrule
    \end{tabular}
  \end{table}

  \metricbox{Template parameterization \texttt{CsrMatrix<IT, NT, OT=IT>} saves \textbf{11.18 GB} of DRAM transfer per SpMV iteration!}
\end{frame}

\begin{frame}[fragile]{SPCraft CSR Matrix Template Architecture}
  \begin{lstlisting}[style=lecturecode]
namespace spcraft {

template <class IT, class NT, class OT = IT>
class CsrMatrix {
public:
  OT* row_ptr = nullptr;  //!< Row pointers (size: m + 1), uses 64-bit OT
  IT* col_id  = nullptr;  //!< Column indices (size: nnz),   uses 32-bit IT
  NT* val     = nullptr;  //!< Numerical values (size: nnz),  uses NT (float/double)

  OT nnz = 0;             //!< Number of nonzeros (wide offset)
  IT m = 0;               //!< Row count (index type)
  IT n = 0;               //!< Column count (index type)

  bool memowned = true;   //!< Memory ownership flag (RAII control)

  // Move-only semantics (prevent accidental gigabyte shallow/deep copies)
  CsrMatrix(const CsrMatrix&) = delete;
  CsrMatrix& operator=(const CsrMatrix&) = delete;
  CsrMatrix(CsrMatrix&&) noexcept;
  CsrMatrix& operator=(CsrMatrix&&) noexcept;
};

} // namespace spcraft
  \end{lstlisting}
\end{frame}

%%%%%%%%%%%%%%%%%%%%%%%%%%%%%%%%%%%%%%%%%%%%%%%%%%%%%%%%%%%%%%%%%%%%%%%%%%%%%%%
\section{03: Resource Ownership, RAII, and Compile-Time Safety}
\sectionfoot{03: Ownership \& Move Semantics}
%%%%%%%%%%%%%%%%%%%%%%%%%%%%%%%%%%%%%%%%%%%%%%%%%%%%%%%%%%%%%%%%%%%%%%%%%%%%%%%

\begin{frame}[fragile]{Compile-Time Error Prevention: The \texttt{e02} Example}
  \begin{columns}[T,onlytextwidth]
    \begin{column}{0.50\textwidth}
      \subhead{Why Copying is Deleted (\texttt{= delete})}
      \begin{lstlisting}[style=lecturecode]
// examples/e02_compiler_error.cpp
#include "spcraft.h"

int main() {
  CsrMatrix<int32_t, float> A;
  // What happens if a user writes:
  // CsrMatrix<int32_t, float> B = A;
}
      \end{lstlisting}
      \begin{alertblock}{Compiler Output}
        \fontsize{6.2}{7.2}\selectfont\texttt{error: call to deleted constructor of 'spcraft::CsrMatrix<int32\_t, float>'}
      \end{alertblock}
      \vspace{0.3em}
      \begin{itemize}\setlength{\itemsep}{0.25em}
        \item Prevents \textbf{accidental double-free crashes}.
        \item Prevents \textbf{silent multi-gigabyte stalls}.
      \end{itemize}
    \end{column}
    \begin{column}{0.46\textwidth}
      \subhead{The Three Explicit Pathways}
      \begin{enumerate}\setlength{\itemsep}{0.6em}
        \item \textbf{Transfer Ownership (Zero Cost)}:
          \begin{lstlisting}[style=lecturecode]
CSR B = std::move(A); // O(1) ptr swap
          \end{lstlisting}
        \item \textbf{Explicit Deep Clone}:
          \begin{lstlisting}[style=lecturecode]
CSR B = A.Clone();    // Intentional alloc
          \end{lstlisting}
        \item \textbf{Non-Owning View (Zero Copy)}:
          \begin{lstlisting}[style=lecturecode]
CSR view(row_ptr, col_id, val,
         nnz, m, n); // memowned=false
          \end{lstlisting}
      \end{enumerate}
    \end{column}
  \end{columns}
\end{frame}

%%%%%%%%%%%%%%%%%%%%%%%%%%%%%%%%%%%%%%%%%%%%%%%%%%%%%%%%%%%%%%%%%%%%%%%%%%%%%%%
\section{04: C++20 Concepts \& Expressive Compile-Time Diagnostics}
\sectionfoot{04: C++20 Concepts}
%%%%%%%%%%%%%%%%%%%%%%%%%%%%%%%%%%%%%%%%%%%%%%%%%%%%%%%%%%%%%%%%%%%%%%%%%%%%%%%

\begin{frame}[fragile]{Transforming Cryptic Compiler Errors into Clear Contracts}
  \begin{columns}[T,onlytextwidth]
    \begin{column}{0.48\textwidth}
      \subhead{C++20 Concepts Definition}
      \begin{lstlisting}[style=lecturecode]
template <typename T>
concept SparseIndex =
  std::integral<T> && (sizeof(T) >= 2);

template <typename T>
concept SparseOffset =
  std::integral<T> && (sizeof(T) >= 4);

template <typename T>
concept SparseValue =
  std::floating_point<T> || std::integral<T>;
      \end{lstlisting}
      \vspace{-0.2cm}
      \begin{itemize}\setlength{\itemsep}{0.3em}
        \item Replaces fragile SFINAE (\texttt{std::enable\_if\_t}) with clean mathematical contracts.
        \item Validates types before compiler attempts instantiation.
      \end{itemize}
    \end{column}
    \begin{column}{0.48\textwidth}
      \subhead{Static Assertions \& Constrained Functions}
      \begin{lstlisting}[style=lecturecode]
template <SparseIndex IT,
          SparseValue NT,
          SparseOffset OT>
requires (sizeof(OT) >= sizeof(IT))
void SpMV(const CsrMatrix<IT, NT, OT>& A,
          const NT* x, NT* y);
      \end{lstlisting}
      \begin{exampleblock}{Human-Readable Diagnostic}
        \fontsize{6.2}{7.2}\selectfont\texttt{error: constraints not satisfied for 'SpMV'\\
        note: 'std::string' does not satisfy 'SparseValue'}
      \end{exampleblock}
    \end{column}
  \end{columns}
\end{frame}

%%%%%%%%%%%%%%%%%%%%%%%%%%%%%%%%%%%%%%%%%%%%%%%%%%%%%%%%%%%%%%%%%%%%%%%%%%%%%%%
\section{05: Static Dispatch \& Branch Elimination with \texttt{if constexpr}}
\sectionfoot{05: \texttt{if constexpr} Static Dispatch}
%%%%%%%%%%%%%%%%%%%%%%%%%%%%%%%%%%%%%%%%%%%%%%%%%%%%%%%%%%%%%%%%%%%%%%%%%%%%%%%

\begin{frame}[fragile]{Eliminating Branches from Innermost SpMV Loops}
  \subhead{Runtime \texttt{if} vs Compile-Time \texttt{if constexpr}}
  \begin{lstlisting}[style=lecturecode]
template <MatrixSymmetry Symmetry = MatrixSymmetry::General, typename IT, typename NT, typename OT>
void FastSpMV(const CsrMatrix<IT, NT, OT>& A, const NT* x, NT* y) {
  if constexpr (Symmetry == MatrixSymmetry::General) {
    // Branch A: Standard CSR SpMV loop (Pure vectorizable streaming loop)
    for (IT i = 0; i < A.m; ++i) {
      NT sum{};
      for (OT k = A.row_ptr[i]; k < A.row_ptr[i + 1]; ++k) sum += A.val[k] * x[A.col_id[k]];
      y[i] += sum;
    }
  } else if constexpr (Symmetry == MatrixSymmetry::Symmetric) {
    // Branch B: Symmetric SpMV (Eliminates storing the lower triangle in memory)
    for (IT i = 0; i < A.m; ++i) {
      for (OT k = A.row_ptr[i]; k < A.row_ptr[i + 1]; ++k) {
        y[i] += A.val[k] * x[A.col_id[k]];
        if (i != A.col_id[k]) y[A.col_id[k]] += A.val[k] * x[i];
      }
    }
  }
}
  \end{lstlisting}
  \metricbox{The unselected branch is completely discarded during compilation. \textbf{Zero runtime branch penalty!}}
\end{frame}

%%%%%%%%%%%%%%%%%%%%%%%%%%%%%%%%%%%%%%%%%%%%%%%%%%%%%%%%%%%%%%%%%%%%%%%%%%%%%%%
\section{06: Expression Templates \& Single-Pass Kernel Fusion}
\sectionfoot{06: Expression Templates}
%%%%%%%%%%%%%%%%%%%%%%%%%%%%%%%%%%%%%%%%%%%%%%%%%%%%%%%%%%%%%%%%%%%%%%%%%%%%%%%

\begin{frame}[fragile]{Expression Templates: Eliminating Intermediate Buffers}
  \begin{columns}[T,onlytextwidth]
    \begin{column}{0.48\textwidth}
      \subhead{The Naive Evaluation Pitfall}
      \begin{equation*}
        \mathbf{y} = \alpha (\mathbf{A} \mathbf{x}) + \beta \mathbf{z}
      \end{equation*}
      \begin{enumerate}\setlength{\itemsep}{0.2em}
        \item \texttt{t1 = A * x} $\implies$ Allocates vector \texttt{t1}
        \item \texttt{t2 = alpha * t1} $\implies$ Allocates vector \texttt{t2}
        \item \texttt{t3 = beta * z} $\implies$ Allocates vector \texttt{t3}
        \item \texttt{y = t2 + t3} $\implies$ 4 memory roundtrips!
      \end{enumerate}
      \begin{alertblock}{Cache Pollution}
        Multiple passes evict cached vectors and overwhelm memory controllers.
      \end{alertblock}
    \end{column}
    \begin{column}{0.48\textwidth}
      \subhead{Lazy Evaluation via AST Trees}
      \begin{center}
        \begin{tikzpicture}[scale=0.85, every node/.style={transform shape}]
          \node[flow] (add) {AddExpr ($+$)};
          \node[data, below left=0.6cm and 0.4cm of add] (scale1) {ScaleExpr ($\alpha \cdot$)};
          \node[data, below right=0.6cm and 0.4cm of add] (scale2) {ScaleExpr ($\beta \cdot$)};
          \node[flow, below=0.5cm of scale1] (spmv) {SpMVExpr ($A \cdot x$)};
          \node[flow, below=0.5cm of scale2] (vec) {VecWrapper ($z$)};

          \draw[arrow] (scale1) -- (add);
          \draw[arrow] (scale2) -- (add);
          \draw[arrow] (spmv) -- (scale1);
          \draw[arrow] (vec) -- (scale2);
        \end{tikzpicture}
      \end{center}
      \vspace{-0.3cm}
      \begin{exampleblock}{Fused Streaming Loop}
        Evaluates $y[i] = \alpha \sum A_{ik} x_k + \beta z_i$ in a \textbf{single streaming pass} directly to DRAM!
      \end{exampleblock}
    \end{column}
  \end{columns}
\end{frame}

%%%%%%%%%%%%%%%%%%%%%%%%%%%%%%%%%%%%%%%%%%%%%%%%%%%%%%%%%%%%%%%%%%%%%%%%%%%%%%%
\section{07: Policy-Based Design for Extensible Sparse Frameworks}
\sectionfoot{07: Policy-Based Design}
%%%%%%%%%%%%%%%%%%%%%%%%%%%%%%%%%%%%%%%%%%%%%%%%%%%%%%%%%%%%%%%%%%%%%%%%%%%%%%%

\begin{frame}[fragile]{Composing Matrix Behavior at Compile Time}
  \subhead{Orthogonal Policies via Templates}
  \begin{lstlisting}[style=lecturecode]
template <typename IT, typename NT, typename OT = IT,
          template <typename, typename, typename> class StoragePolicy = OwningStoragePolicy,
          typename IndexPolicy = ZeroBasedPolicy>
class PolicySparseMatrix : public StoragePolicy<IT, NT, OT>, public IndexPolicy {
  // Composes behaviors without any virtual table or pointer overhead!
};

// 1. Standard C++ CSR (0-indexed, heap-owning)
using StandardCsr = PolicySparseMatrix<int32_t, double, int32_t, OwningStoragePolicy, ZeroBasedPolicy>;

// 2. MatrixMarket / Fortran / MATLAB CSR (1-indexed, heap-owning)
using FortranCsr  = PolicySparseMatrix<int32_t, double, int32_t, OwningStoragePolicy, OneBasedPolicy>;

// 3. Zero-copy GPU / External View (0-indexed, non-owning)
using ViewCsr     = PolicySparseMatrix<int32_t, double, int64_t, ViewStoragePolicy, ZeroBasedPolicy>;
  \end{lstlisting}
  \metricbox{All policy member methods (\texttt{ToInternal}, \texttt{AllocateStorage}) are \textbf{100\% inlined}.}
\end{frame}

%%%%%%%%%%%%%%%%%%%%%%%%%%%%%%%%%%%%%%%%%%%%%%%%%%%%%%%%%%%%%%%%%%%%%%%%%%%%%%%
\section{Summary \& Engineering Best Practices}
\sectionfoot{Summary}
%%%%%%%%%%%%%%%%%%%%%%%%%%%%%%%%%%%%%%%%%%%%%%%%%%%%%%%%%%%%%%%%%%%%%%%%%%%%%%%

\begin{frame}{Modern C++ Metaprogramming Checklist for Sparse HPC}
  \begin{columns}[T,onlytextwidth]
    \begin{column}{0.48\textwidth}
      \subhead{Rules to Follow}
      \begin{itemize}\setlength{\itemsep}{0.5em}
        \item[\checkmark] \textbf{Decouple IT, NT, OT}: Allow 64-bit offsets with 32-bit indices for hyper-sparse and large graphs.
        \item[\checkmark] \textbf{Delete Copy Operations}: Force users to consciously choose \texttt{std::move}, \texttt{.Clone()}, or non-owning views.
        \item[\checkmark] \textbf{Constrain with C++20 Concepts}: Provide clean, actionable error messages instead of cryptic template dumps.
      \end{itemize}
    \end{column}
    \begin{column}{0.48\textwidth}
      \subhead{Performance Maximizers}
      \begin{itemize}\setlength{\itemsep}{0.5em}
        \item[\checkmark] \textbf{Use \texttt{if constexpr}}: Eliminate branch conditions from tight computational loops.
        \item[\checkmark] \textbf{Leverage Inlining (\texttt{-inl.h})}: Keep template definitions accessible to the compiler for SIMD auto-vectorization.
        \item[\checkmark] \textbf{Apply Expression Templates}: Fuse composite matrix/vector expressions to reduce DRAM memory roundtrips.
      \end{itemize}
    \end{column}
  \end{columns}
  \vfill
  \metricbox{All executable examples available in: \texttt{examples/metaprogramming/}}
\end{frame}

\end{document}
